\section{KL 散度与 Fisher 度量的局部关系}
\label{sec:12-Real-Analysis-Support/09-Information-Geometry/03}

调过策略优化算法的人都见过那行配置：\texttt{target\_kl = 0.01}。新旧策略之间的 KL 散度一旦超过它，这一轮更新就被截断或回退。奇怪的地方在于，人们想控制的明明是``别把策略改得太狠''，落到代码里却成了一个信息论的量。为什么不直接限制参数的移动幅度？答案藏在 KL 散度的一个局部性质里：在参数邻域内，它是一个二次型，而这个二次型的矩阵正是 Fisher 信息。

\subsection{一阶项因 score 期望为零而消失}

设 \(p_\theta\) 光滑地依赖参数，考察小位移 \(\Delta\) 带来的 KL：

\[
\KL\bigl(p_\theta\,\|\,p_{\theta+\Delta}\bigr)
=\E_\theta\!\left[\log p_\theta(X)-\log p_{\theta+\Delta}(X)\right].
\]

记 score 为 \(s(x)=\nabla_\theta\log p_\theta(x)\)，把被积项对 \(\Delta\) 展开到二阶：

\[
\log p_{\theta+\Delta}(x)=\log p_\theta(x)+s(x)^{\trans}\Delta
+\tfrac12\,\Delta^{\trans}\nabla_\theta^2\log p_\theta(x)\,\Delta
+o\bigl(\norm{\Delta}^2\bigr).
\]

代回去取期望，得到

\[
\begin{aligned}
\KL\bigl(p_\theta\,\|\,p_{\theta+\Delta}\bigr)
&=-\E_\theta[s(X)]^{\trans}\Delta
-\tfrac12\,\Delta^{\trans}\E_\theta\bigl[\nabla_\theta^2\log p_\theta(X)\bigr]\Delta
+o\bigl(\norm{\Delta}^2\bigr)\\[2pt]
&=\tfrac12\,\Delta^{\trans}I(\theta)\,\Delta+o\bigl(\norm{\Delta}^2\bigr).
\end{aligned}
\]

两步替换值得单独点出。一阶项整体消失，因为 \(\E_\theta[s(X)]=0\)：对恒等式 \(\int p_\theta\,d\mu=1\) 两边求导，导数穿过积分号后得到 \(\int\nabla_\theta p_\theta\,d\mu=0\)，而 \(\nabla_\theta p_\theta=p_\theta\,s\)。二阶项用的是信息等式 \(I(\theta)=\E_\theta[ss^{\trans}]=-\E_\theta[\nabla_\theta^2\log p_\theta]\)，它由同一个恒等式再求一次导得到，推导见\hyperref[sec:11-Mathematical-Statistics/03-Sufficient-Statistics-and-Information/02]{``Score 与 Fisher 信息衡量局部可辨认性''}。

条件对账：两次``求导穿过积分号''都需要支撑集不随 \(\theta\) 改变、并且有可积的控制函数，条件见\hyperref[sec:12-Real-Analysis-Support/04-Limits-and-Integration/03]{``求导穿过积分号需要控制差商''}。\(\Uniform(0,\theta)\) 这样的模型不满足前一条——支撑集端点就是参数本身，它的 KL 在局部不是二次的，Fisher 信息的通常公式也失效。另外还要求 \(I(\theta)\) 有限，且 \(\log p_\theta\) 关于 \(\theta\) 二阶连续可微。

\subsection{等值线的形状由 Fisher 的特征分解决定}

把上式反过来读：在 \(\theta_0\) 的一个小邻域里，\(\KL(p_{\theta_0}\|p_\theta)\approx\tfrac12(\theta-\theta_0)^{\trans}I(\theta_0)(\theta-\theta_0)\) 是一个正定二次型，它的等值集是椭球。

\begin{center}
\begin{tikzpicture}[>=stealth]
\begin{scope}[shift={(3.4,2.2)}]
\draw[gray!60,dashed] (0,0) circle (1.0);
\foreach \s in {1,1.414,2}
  \draw[rotate around={30:(0,0)}] (0,0) ellipse ({1.0*\s} and {0.42*\s});
\draw[->,thick,rotate around={30:(0,0)}] (0,0) -- (2.0,0);
\draw[->,thick] (0,0) -- (0.42,-0.727);
\fill (0,0) circle (0.06);
\node[fill=white,inner sep=1pt,left] at (-0.10,-0.16) {\(\theta_0\)};
\node[right] at (1.80,1.06) {\(I\) 的小特征值方向};
\node[below right] at (0.46,-0.80) {大特征值方向};
\draw[gray!70] (-1.45,1.32) -- (-0.50,0.866);
\node[gray!80,left] at (-1.45,1.32) {欧氏圆};
\draw[gray!70] (-2.05,-1.36) -- (-1.45,-1.17);
\node[left] at (-2.10,-1.38) {\(\KL=\varepsilon,\,2\varepsilon,\,4\varepsilon\)};
\end{scope}
\end{tikzpicture}
\end{center}

\emph{KL 在 \(\theta_0\) 附近的等值线是椭圆，主轴方向与长短由 Fisher 矩阵的特征分解给出；参数空间的欧氏圆（虚线）与它一般不重合。}

把 \(I(\theta_0)=U\Lambda U^{\trans}\) 特征分解，沿第 \(i\) 个特征方向走到 \(\KL=\varepsilon\) 需要的参数位移是 \(\sqrt{2\varepsilon/\lambda_i}\)。特征值大的方向上分布敏感，允许的参数位移小；特征值小的方向上分布迟钝，同样的 KL 预算可以换来很长的一步。椭圆被拉长的那一头，正是模型``不在乎''的方向。上一节信任域图里那个 Fisher 椭圆，与这里的 KL 等值线是同一个东西——所以约束``分布改变不超过 \(\varepsilon\)''与约束``二次型不超过 \(2\varepsilon\)''在局部是一回事。

椭圆的形状还随位置改变。走到参数空间的另一处，\(I(\theta)\) 变了，椭圆的朝向和长短跟着变。固定半径的欧氏球做不到这种自适应，这也是限制参数范数与限制 KL 的分歧所在。

\subsection{不对称性要到三阶才现身}

KL 散度两个方向不相等，而二次型是对称的，这两件事需要对上账。把展开点固定在 \(\theta\)，反方向的 KL 给出

\[
\KL\bigl(p_{\theta+\Delta}\,\|\,p_\theta\bigr)
=\tfrac12\,\Delta^{\trans}I(\theta)\,\Delta+o\bigl(\norm{\Delta}^2\bigr),
\]

与正方向的二阶项完全相同。差别从三阶项开始出现，由 \(I(\theta+\Delta)\) 与 \(I(\theta)\) 的偏离所贡献。于是有限位移下前向与反向 KL 的著名分歧——一个倾向覆盖住目标分布的全部支撑，一个倾向缩进目标分布的高密度区，见\hyperref[sec:07-Information-Theory/03-Divergences/02]{``Gibbs 不等式与 KL 的方向性''}——在无穷小极限下消失干净，两个方向诱导出同一个黎曼度量。度量的对称性由此而来，并不是额外加上去的假设。

顺带说清楚一层关系：KL 本身不是距离的平方，它没有对称性也没有三角不等式。局部二次近似为它配上了一个合格的度量，而这个度量只在无穷小的意义上代表它。

\subsection{信任域方法约束的正是这个椭圆}

回到开头的 \texttt{target\_kl}。TRPO 把一步更新写成

\[
\max_{\theta}\ \widehat L(\theta)
\qquad\text{s.t.}\qquad
\E_{s\sim d}\bigl[\KL(\pi_{\theta_{\rm old}}(\cdot\given s)\,\|\,\pi_\theta(\cdot\given s))\bigr]\le\delta .
\]

把目标线性化、把约束按上面的展开换成二次型，这个问题的解就是自然梯度方向乘一个由 \(\delta\) 决定的步长：

\[
\Delta^{\star}=\sqrt{\frac{2\delta}{\nabla \widehat L^{\trans}I^{-1}\nabla \widehat L}}\;I^{-1}\nabla \widehat L .
\]

上一节从信任域推出自然梯度时用的是同一套代数，只是那里的椭圆来自度量的定义，这里的椭圆来自 KL 的二阶展开——两条路走到同一个式子，是因为这两个椭圆本来就是一个。TRPO 的实现用共轭梯度求解 \(I^{-1}\nabla\widehat L\)，避开显式构造 \(I\)。

PPO 放弃了显式约束，改用截断的重要性比值或一个 KL 惩罚项，它保留``限制分布位移''的意图而牺牲了几何上的精确性。理解了椭圆的图像，几个工程现象就不再费解：\(\delta\) 应当按分布位移而不是按参数尺度来设，因此换一种网络参数化不必重调；策略接近确定性时 \(I\) 的某些特征值急剧变大，同样的 \(\delta\) 只允许极小的参数移动，训练看上去``卡住''；而在 \(I\) 近乎奇异的方向上，KL 约束几乎不起作用，需要阻尼来兜底。

\subsection{局部成立，不等于全局可用}

二阶展开只在小邻域内有效。位移一大，KL 的真实值与二次型可以差出数量级，尤其在分布支撑几乎错开的地方——此时 KL 趋向无穷，而二次型仍是有限数。信任域算法里那道回退检查（算出候选步之后重新计算真实 KL，超了就缩步长）针对的就是这种失配。

另一处边界是退化。过参数化模型的 \(I(\theta)\) 有大量零特征值，椭圆在那些方向上退化成无界的柱体，二次型不再定义距离。实践中的阻尼把 \(I\) 换成 \(I+\lambda\mathrm{Id}\)，等价于给椭圆强行加上一个欧氏的下界。

这三节走完，三个原本分属不同学科的量接在了一处：信息论用 KL 度量``用错模型要多付多少''，见\hyperref[sec:07-Information-Theory/03-Divergences/01]{``错误模型的代价：从交叉熵到 KL 散度''}；统计学用 Fisher 信息刻画``参数有多容易辨认''；优化用度量决定``一步该往哪儿迈''。局部展开表明它们说的是同一件事的三个侧面，而这正是信息几何在参数化、坐标选择与更新方向成为问题时能给出统一答案的原因。
